% ==============================================================
%  Chapter 7: AI for Managers
% ==============================================================
\chapter{AI for Managers}
\label{ch:managers}

Management work is largely information work: synthesising data, communicating decisions,
planning, resolving ambiguity, and aligning people.
AI is a capable assistant for all of these --- but it works for you, not instead of you.
The judgment, relationship, and accountability remain entirely yours.

\section{Strategic Planning and Decision Making}
\label{sec:strategy}

AI is a strong strategy co-pilot.
It generates scenario options, pressure-tests reasoning, surfaces unconsidered risks,
and structures frameworks.
It cannot know your organisation's culture, political dynamics, or recent history unless
you explain them.

\textbf{Strategic planning prompts:}
\begin{itemize}
  \item \textit{``We are a 40-person B2B SaaS company. Our growth has slowed from 80\% to 25\% year-over-year. Generate three strategic responses with tradeoffs.''}
  \item \textit{``What are the top 5 risks in this plan, and for each, what is a mitigation? [paste plan]''}
  \item \textit{``What assumptions is this strategy relying on? Which are most likely to be wrong?''}
  \item \textit{``Steelman the case against this decision. What would a smart critic say?''}
\end{itemize}

\begin{tipbox}
Explicitly ask AI to challenge your thinking.
\textit{``What would a smart critic say is wrong with this plan?''}
AI defaults to affirming your framing unless you actively ask for pushback.
\end{tipbox}

\section{Communication and Reporting}
\label{sec:communication}

Communication is where managers extract the fastest, most reliable value from AI.
Writing takes time, and AI writes well across formats and audiences.

\textbf{Communication prompts by format:}
\begin{itemize}
  \item \textbf{Executive summary:}
    \textit{``Rewrite this technical status update as a 3-paragraph executive summary for a non-technical board. Lead with business impact, not technical details.''}
  \item \textbf{Status report:}
    \textit{``Summarise these project notes as a weekly status report: RAG status, key achievements, blockers, and next week's priorities.''}
  \item \textbf{Difficult message:}
    \textit{``Help me communicate a 2-week schedule slip to stakeholders. Keep it factual, own accountability, and focus on the recovery plan.''}
  \item \textbf{Meeting prep:}
    \textit{``I have a 30-minute meeting with a sceptical VP to pitch this proposal. What objections should I prepare for and how do I address each?''}
\end{itemize}

\section{Team Productivity and Governance}
\label{sec:team-productivity}

\textbf{AI adoption for your team:}
The most important thing you can do as a manager is set a clear, simple AI policy.
Teams with no guidance default to either over-using AI (producing low-quality work fast)
or avoiding it (leaving productivity gains on the table).

A one-page policy should cover:
\begin{enumerate}
  \item \textbf{Approved tools} --- which tools are sanctioned for what categories of work
  \item \textbf{Data rules} --- what data may never be entered into any AI tool (PII, trade
        secrets, financial results, customer names in non-enterprise accounts)
  \item \textbf{Review requirements} --- what categories of AI output require human review
        before use (customer communications, code going to production, legal or financial
        content)
  \item \textbf{Attribution} --- whether AI assistance should be disclosed internally
\end{enumerate}

\begin{warnbox}
Data privacy is the highest-risk area for managers.
An employee pasting a customer's personal data or a commercially sensitive strategy
into a public AI tool is a potential data breach.
This is a training and policy problem, not a technical problem.
\end{warnbox}

\textbf{Measuring AI productivity impact:}
Measure outcomes, not AI usage.
Useful metrics: cycle time, document turnaround time, defect rate, meeting action-item
completion rate.
Compare before and after a 6--8 week adoption period.
Avoid the trap of measuring prompts sent or AI tool logins.

\section{Hiring and Team Composition}
\label{sec:hiring}

AI shifts the value of different skills.
Roles involving high-volume structured output (first-draft writing, standard code
generation, data formatting) are being partially automated.
Roles requiring judgment, context sensitivity, relationship management, and evaluation
of AI output are becoming more valuable.

\textbf{Adjust your hiring criteria:}
\begin{itemize}
  \item Prioritise judgment and critical thinking over syntax or template knowledge
  \item Ask candidates how they use AI in their work and how they verify its output
  \item Be cautious about hiring someone who cannot evaluate AI output in their domain
  \item Junior roles change more than senior roles --- plan for more coaching
\end{itemize}

\section{AI Governance at the Manager Level}
\label{sec:governance}

Every manager needs to answer these questions for their team:

\begin{enumerate}
  \item \textbf{Which tools are approved?} Maintain a short list and update it quarterly.
  \item \textbf{What data is off-limits?} Write this clearly and train the team on it.
  \item \textbf{What output needs human review?} Define this per output category.
  \item \textbf{Who is accountable when AI is wrong?} Always a human --- make this explicit.
  \item \textbf{How do we handle an AI-related incident?} Have a response process ready.
\end{enumerate}

\begin{infobox}[Accountability Principle]
The accountability for any AI output used in your organisation belongs entirely to
the human who used it.
\textit{``The AI said so''} is not a defence --- it is evidence of inadequate review.
Make this unambiguous with your team.
\end{infobox}
